% Options for packages loaded elsewhere
\PassOptionsToPackage{unicode}{hyperref}
\PassOptionsToPackage{hyphens}{url}
\PassOptionsToPackage{dvipsnames,svgnames,x11names}{xcolor}
%
\documentclass[
  letterpaper,
  DIV=11,
  numbers=noendperiod]{scrartcl}

\usepackage{amsmath,amssymb}
\usepackage{iftex}
\ifPDFTeX
  \usepackage[T1]{fontenc}
  \usepackage[utf8]{inputenc}
  \usepackage{textcomp} % provide euro and other symbols
\else % if luatex or xetex
  \usepackage{unicode-math}
  \defaultfontfeatures{Scale=MatchLowercase}
  \defaultfontfeatures[\rmfamily]{Ligatures=TeX,Scale=1}
\fi
\usepackage{lmodern}
\ifPDFTeX\else  
    % xetex/luatex font selection
\fi
% Use upquote if available, for straight quotes in verbatim environments
\IfFileExists{upquote.sty}{\usepackage{upquote}}{}
\IfFileExists{microtype.sty}{% use microtype if available
  \usepackage[]{microtype}
  \UseMicrotypeSet[protrusion]{basicmath} % disable protrusion for tt fonts
}{}
\makeatletter
\@ifundefined{KOMAClassName}{% if non-KOMA class
  \IfFileExists{parskip.sty}{%
    \usepackage{parskip}
  }{% else
    \setlength{\parindent}{0pt}
    \setlength{\parskip}{6pt plus 2pt minus 1pt}}
}{% if KOMA class
  \KOMAoptions{parskip=half}}
\makeatother
\usepackage{xcolor}
\setlength{\emergencystretch}{3em} % prevent overfull lines
\setcounter{secnumdepth}{-\maxdimen} % remove section numbering
% Make \paragraph and \subparagraph free-standing
\makeatletter
\ifx\paragraph\undefined\else
  \let\oldparagraph\paragraph
  \renewcommand{\paragraph}{
    \@ifstar
      \xxxParagraphStar
      \xxxParagraphNoStar
  }
  \newcommand{\xxxParagraphStar}[1]{\oldparagraph*{#1}\mbox{}}
  \newcommand{\xxxParagraphNoStar}[1]{\oldparagraph{#1}\mbox{}}
\fi
\ifx\subparagraph\undefined\else
  \let\oldsubparagraph\subparagraph
  \renewcommand{\subparagraph}{
    \@ifstar
      \xxxSubParagraphStar
      \xxxSubParagraphNoStar
  }
  \newcommand{\xxxSubParagraphStar}[1]{\oldsubparagraph*{#1}\mbox{}}
  \newcommand{\xxxSubParagraphNoStar}[1]{\oldsubparagraph{#1}\mbox{}}
\fi
\makeatother


\providecommand{\tightlist}{%
  \setlength{\itemsep}{0pt}\setlength{\parskip}{0pt}}\usepackage{longtable,booktabs,array}
\usepackage{calc} % for calculating minipage widths
% Correct order of tables after \paragraph or \subparagraph
\usepackage{etoolbox}
\makeatletter
\patchcmd\longtable{\par}{\if@noskipsec\mbox{}\fi\par}{}{}
\makeatother
% Allow footnotes in longtable head/foot
\IfFileExists{footnotehyper.sty}{\usepackage{footnotehyper}}{\usepackage{footnote}}
\makesavenoteenv{longtable}
\usepackage{graphicx}
\makeatletter
\def\maxwidth{\ifdim\Gin@nat@width>\linewidth\linewidth\else\Gin@nat@width\fi}
\def\maxheight{\ifdim\Gin@nat@height>\textheight\textheight\else\Gin@nat@height\fi}
\makeatother
% Scale images if necessary, so that they will not overflow the page
% margins by default, and it is still possible to overwrite the defaults
% using explicit options in \includegraphics[width, height, ...]{}
\setkeys{Gin}{width=\maxwidth,height=\maxheight,keepaspectratio}
% Set default figure placement to htbp
\makeatletter
\def\fps@figure{htbp}
\makeatother

\KOMAoption{captions}{tableheading}
\makeatletter
\@ifpackageloaded{caption}{}{\usepackage{caption}}
\AtBeginDocument{%
\ifdefined\contentsname
  \renewcommand*\contentsname{Table of contents}
\else
  \newcommand\contentsname{Table of contents}
\fi
\ifdefined\listfigurename
  \renewcommand*\listfigurename{List of Figures}
\else
  \newcommand\listfigurename{List of Figures}
\fi
\ifdefined\listtablename
  \renewcommand*\listtablename{List of Tables}
\else
  \newcommand\listtablename{List of Tables}
\fi
\ifdefined\figurename
  \renewcommand*\figurename{Figure}
\else
  \newcommand\figurename{Figure}
\fi
\ifdefined\tablename
  \renewcommand*\tablename{Table}
\else
  \newcommand\tablename{Table}
\fi
}
\@ifpackageloaded{float}{}{\usepackage{float}}
\floatstyle{ruled}
\@ifundefined{c@chapter}{\newfloat{codelisting}{h}{lop}}{\newfloat{codelisting}{h}{lop}[chapter]}
\floatname{codelisting}{Listing}
\newcommand*\listoflistings{\listof{codelisting}{List of Listings}}
\makeatother
\makeatletter
\makeatother
\makeatletter
\@ifpackageloaded{caption}{}{\usepackage{caption}}
\@ifpackageloaded{subcaption}{}{\usepackage{subcaption}}
\makeatother

\ifLuaTeX
  \usepackage{selnolig}  % disable illegal ligatures
\fi
\usepackage{bookmark}

\IfFileExists{xurl.sty}{\usepackage{xurl}}{} % add URL line breaks if available
\urlstyle{same} % disable monospaced font for URLs
\hypersetup{
  pdftitle={The Spatial Geography of Roman Warfare},
  colorlinks=true,
  linkcolor={blue},
  filecolor={Maroon},
  citecolor={Blue},
  urlcolor={Blue},
  pdfcreator={LaTeX via pandoc}}


\title{The Spatial Geography of Roman Warfare}
\usepackage{etoolbox}
\makeatletter
\providecommand{\subtitle}[1]{% add subtitle to \maketitle
  \apptocmd{\@title}{\par {\large #1 \par}}{}{}
}
\makeatother
\subtitle{Battle Locations, Territorial Boundaries, Settlements, and
Military Infrastructure}
\author{John Grace}
\date{2026-11-22}

\begin{document}
\maketitle
\begin{abstract}
This study examines the spatial distribution of Roman battles and
investigates the relationship between military conflict, settlement
patterns, territorial boundaries, and specialized infrastructure across
the Roman world. Using georeferenced battle locations, historical
territorial boundary reconstructions, and settlement records from the
Pleiades Gazetteer, the study applies spatial statistical methods and
Monte Carlo simulations to evaluate whether battle locations exhibit
patterns beyond those expected under alternative spatial distributions.
Particular attention is given to differences between external conflicts
and internal Roman civil conflicts. The analysis considers spatial
randomness, settlement intensity, proximity to territorial boundaries,
and proximity to military and civic infrastructure.
\end{abstract}


\section{Introduction}\label{introduction}

Warfare was a fundamental component of the political, territorial, and
administrative development of the Roman world. Military engagements
occurred across a geographically extensive territory, ranging from the
Italian peninsula and western Mediterranean to the frontiers of northern
Europe, North Africa, and the eastern Mediterranean. Although the
historical circumstances of individual battles have received
considerable scholarly attention, their collective geographic
distribution has been understudied. Utilizing geographical statistical
modeling offers an avenue to build upon past stragetic history on the
Roman Empire/.

Battle locations were not necessarily distributed randomly across the
landscape. Military engagements may have been influenced by the location
of established settlements, the boundaries of Roman territorial control,
and the presence of specialized infrastructure such as forts. Internal
conflicts may exhibit different geographic characteristics from
conflicts involving external opponents.

By combining georeferenced battle records with reconstructed historical
boundaries and settlement data, it is possible to evaluate whether
observed battle locations exhibit spatial patterns consistent with
particular geographic features.

This study investigates the spatial distribution of Roman battles
through a series of spatial statistical tests and Monte Carlo
simulations. Rather than assuming that a single geographic factor
explains the location of military engagements, the analysis compares
several alternative spatial models.

The primary objectives are:

\begin{enumerate}
\def\labelenumi{\arabic{enumi}.}
\tightlist
\item
  Determine whether Roman battle locations exhibit significant spatial
  clustering relative to complete spatial randomness.
\item
  Evaluate whether the distribution of battles is associated with the
  historical distribution of settlements.
\item
  Investigate whether battles occurred disproportionately close to the
  territorial boundaries of Roman control.
\item
  Examine whether external conflicts and civil conflicts exhibit
  different spatial relationships with military and civic
  infrastructure.
\end{enumerate}

By examining these relationships together, this study seeks to better
understand the geographic characteristics of Roman warfare and the
relationship between military activity and the historical Roman
landscape.

\section{Data and Study Area}\label{data-and-study-area}

\subsection{Roman Battle Locations}\label{roman-battle-locations}

The primary dataset consists of georeferenced battle records associated
with Roman military history.

Battle records were collected using Wikidata and subsequently reviewed
to determine whether Rome was involved in each engagement. The resulting
dataset includes battle identifiers, names, dates, classifications, and
associated geographic locations.

Each battle is assigned to one of two classifications:

\begin{itemize}
\tightlist
\item
  \textbf{External:} Conflicts involving Rome and external opponents.
\item
  \textbf{Civil Conflict:} Conflicts involving internal Roman political
  or military factions.
\end{itemize}

Battle locations are represented using geographic coordinates and
converted into spatial point geometries for subsequent analysis. Because
individual battles may have multiple associated location records, the
analysis distinguishes between battle-location records and unique
battles. Analyses using multiple location records retain the associated
geographic observations, while the later classification-based nearest
neighbor analysis constructs one representative point per unique battle.

\subsection{Historical Roman Territorial
Boundaries}\label{historical-roman-territorial-boundaries}

Historical Roman territorial boundaries were collected to represent the
changing geographic extent of Roman control across nine historical
periods. Boundary reconstructions were obtained from two online
geographic data sources: the \emph{Roman Empire Cultural Shapefiles}
ArcGIS Feature Service and the \emph{Roman Empire Map} project.

The ArcGIS Feature Service provided six historical boundary layers. The
available layers were first inspected through the service's metadata to
identify relevant historical reconstructions and their associated
geographic attributes. The selected layers were downloaded as GeoJSON
files using the ArcGIS REST API, with their output coordinate reference
system specified as WGS 84 (EPSG:4326).

Three additional Republican period boundary reconstructions were
obtained from the Roman Empire Map project's GeoJSON dataset. The source
dataset contained multiple historical boundary features, which were
filtered by date and by the \texttt{part} attribute. Only features
representing the complete territorial extent (\texttt{part\ =\ "whole"})
were retained for the years 264 BC, 133 BC, and 44 BC. These selected
boundaries were then converted to EPSG:4326 and saved as individual
GeoJSON files.

The resulting collection contains nine historical boundary snapshots:

\begin{longtable}[]{@{}lrl@{}}
\toprule\noalign{}
Historical snapshot & Year & Data source \\
\midrule\noalign{}
\endhead
\bottomrule\noalign{}
\endlastfoot
264 BC & -264 & Roman Empire Map \\
133 BC & -133 & Roman Empire Map \\
60 BC & -60 & ArcGIS Feature Service \\
44 BC & -44 & Roman Empire Map \\
14 AD & 14 & ArcGIS Feature Service \\
69 AD & 69 & ArcGIS Feature Service \\
117 AD & 117 & ArcGIS Feature Service \\
200 AD & 200 & ArcGIS Feature Service \\
314 AD & 314 & ArcGIS Feature Service \\
\end{longtable}

After collection, the boundary files were loaded into GeoPandas for
inspection. This included reviewing their feature counts, coordinate
reference systems, geometry types, geographic extents, and geometry
validity. The individual reconstructions were also mapped separately and
overlaid to examine how the geographic extent of Roman territory changed
across the selected periods.

These boundary datasets provide historical approximations of Roman
territorial extent rather than continuous representations of Roman
expansion and contraction. Because the reconstructions are available
only for selected historical dates, they serve as temporal reference
points for the spatial analysis rather than exact representations of
Roman territorial control at every battle date.

\subsection{Settlement Data}\label{settlement-data}

Settlement information is obtained from the Pleiades Gazetteer, a
historical geographic resource documenting locations associated with the
ancient world. The settlement dataset includes geographic coordinates,
feature classifications, and temporal attributes. Two temporal
attributes are particularly important:

\begin{itemize}
\tightlist
\item
  \texttt{MINDATE}: Earliest associated date.
\item
  \texttt{MAXDATE}: Latest associated date.
\end{itemize}

These attributes are used to construct historical settlement datasets
corresponding to the nine Roman territorial snapshots. A settlement is
included in a historical period when its recorded date range contains
the year of that snapshot. The settlement dataset also includes
specialized feature classifications, allowing the analysis to
distinguish between general settlements and particular types of
infrastructure. These include ports, temples, forts, amphitheaters, and
theaters.

\subsection{Geographic Study Area}\label{geographic-study-area}

The geographic study area is defined using the Roman territorial
boundary reconstruction from 314 AD. The boundary geometry is combined
and projected into an coordinate reference system centered approximately
on the Mediterranean region. A 500-kilometer buffer is then constructed
around the territorial boundary. This buffer allows for Roman campaigns
deep into territory that they did not control, such as campaigns fought
agaisnt Parthia or Sassanid Iran in territory that they would never
control, but also for earlier periods allow for battles in Spain and
North Africa during the punic wars that would at later periods be
incorporated into the Roman Empire. Settlement records are restricted to
this study area before the primary spatial analyses.

\section{Methodology}\label{methodology}

The analysis is organized as a sequence of spatial investigations
designed to develop increasingly specific explanations for the
geographic distribution of Roman battles. Rather than treating battle
locations as an isolated spatial phenomenon, the analysis progressively
incorporates historical and geographic context to examine the factors
associated with their distribution.

The first stage establishes the degree of spatial clustering among
recorded battles. The second examines whether the distribution of
settlements provides an explanation for this clustering. The third
introduces Roman territorial boundaries to investigate the relationship
between military activity and the geographic limits of Roman control.
Finally, the analysis considers whether different categories of conflict
exhibit distinct relationships with specialized military and civic
infrastructure.

\subsection{Spatial Data Preparation}\label{spatial-data-preparation}

The battle dataset contains geographic records associated with
individual battles. Records with missing or invalid coordinates are
excluded before spatial analysis, and the remaining coordinates are
converted into point geometries. Settlement locations are similarly
restricted to valid point geometries within the defined study area.

All spatial datasets are transformed into a common projected coordinate
reference system to support distance calculations in meters. The study
area is defined using the historical Roman territorial boundary for 314
AD and a 500 km surrounding buffer.

The unit of analysis varies across the investigations. The initial
spatial randomness, settlement, and territorial-boundary analyses use
battle-location records. The final classification-based infrastructure
analysis instead consolidates records belonging to the same battle into
a single representative location. For battles with multiple location
records, the geometries are combined and their centroid is used as the
representative point. This produces a dataset of unique battles for
comparing the spatial clustering of external and civil conflicts.

\subsection{Monte Carlo Simulation
Framework}\label{monte-carlo-simulation-framework}

The primary statistical approach uses Monte Carlo simulations to compare
observed battle distributions against alternative spatial models.

For each model, the observed spatial statistic is compared with the
empirical distribution generated by its corresponding simulations. The
expected statistic is calculated as the mean of the simulated values,
while the 2.5th and 97.5th percentiles provide an empirical interval for
evaluating the observed result. The observed-to-expected ratio is
calculated by dividing the observed statistic by the simulated mean. For
distance-based statistics, ratios below 1 indicate that the observed
distances are shorter than expected, while ratios above 1 indicate that
they are longer.

The direction of the Monte Carlo test depends on the research question.
Lower-tail tests are used when evaluating whether observed distances are
unusually short, such as in the settlement-feature proximity analysis.
Two-sided tests are used for the classification-based infrastructure
models, where unusually short and unusually long nearest-neighbor
distances are both relevant.

For a lower-tail test, the empirical Monte Carlo p-value is calculated
as:

{[} p = \frac{b+1}{R+1} {]}

where (b) is the number of simulated statistics less than or equal to
the observed statistic and (R) is the number of simulations. For
two-sided tests, the p-value is calculated by doubling the smaller of
the two empirical tail probabilities, with a maximum value of 1.

With 999 simulations, the minimum attainable p-value is 0.001 for a
one-sided test and 0.002 for the two-sided procedure. P-values are
interpreted as evidence of differences between the observed and
simulated spatial statistics rather than as measures of the probability
that a particular historical explanation is true.

\subsection{Spatial Clustering}\label{spatial-clustering}

The initial analysis evaluates whether Roman battle locations exhibit
greater spatial clustering than would be expected under complete spatial
randomness (CSR). Mean nearest-neighbor distance is used as the measure
of spatial concentration. For each observed battle location, a k-d tree
is used to identify the distance to its nearest neighboring battle
location. The mean of these distances provides the observed clustering
statistic.

The CSR null model generates simulated battle locations by uniformly
sampling points within the study area. Each simulation preserves the
number of observed battle-location records. The mean nearest-neighbor
distance is calculated for each simulated distribution and compared with
the observed statistic. Smaller mean distances indicate greater spatial
clustering.

\subsection{Settlement Geography}\label{settlement-geography}

The settlement-based model evaluates whether the geographic distribution
of documented settlements helps explain the spatial concentration of
Roman battles. Because reliable population-density estimates are
difficult to establish for the ancient world, documented settlement
locations are used as an imperfect proxy for the spatial distribution of
human activity.

A Gaussian kernel density estimate (KDE) is fitted to the filtered
Pleiades settlement locations using a bandwidth of 50 km. The resulting
intensity surface is evaluated across a 20 km sampling grid covering the
study area. Grid-cell intensity estimates are normalized into sampling
probabilities, with locations near concentrations of documented
settlements receiving greater probability of selection.

Monte Carlo simulations generate synthetic battle locations by sampling
grid cells according to these probabilities. Each simulation preserves
the number of observed battle locations and is used to calculate an
expected mean nearest-neighbor distance. This provides a
settlement-conditioned baseline for evaluating whether the observed
clustering of battles can be explained by the distribution of known
settlements.

For the historical boundary analysis, separate settlement intensity
surfaces are also constructed for each of the nine historical boundary
snapshots. Settlement records are assigned to a snapshot using their
Pleiades temporal attributes, retaining locations whose recorded date
ranges include the snapshot year.

The settlement-based model is also applied separately to external
battles and civil conflicts. These classification-specific simulations
preserve the number of observed battle locations in each group and
provide separate expected nearest-neighbor distances for comparison with
the corresponding observed distributions. The combined results provide a
general settlement-based baseline for evaluating the
infrastructure-weighted models.

\subsection{Territorial Boundaries}\label{territorial-boundaries}

The third stage introduces historical territorial boundaries to examine
the relationship between Roman military activity and the geographic
extent of Roman control. Battle locations are compared with the nearest
historical territorial perimeter corresponding to their assigned
chronological period. Their observed distances from these boundaries are
evaluated against distances generated by settlement-based simulations.
This approach examines whether the proximity of battles to territorial
boundaries can be explained by the underlying distribution of
settlements, or whether battles exhibit an additional concentration near
the boundaries of Roman territory.

Territorial boundary association is examined by calculating the shortest
distance between each observed battle location and the corresponding
historical Roman territorial perimeter.

For each battle, the applicable perimeter is determined using its
assigned historical boundary snapshot. Each battle is assigned to the
closest available historical boundary snapshot based on its date. The
analysis uses nine historical snapshots, ranging from 264 BC to 314 AD.
This approach allows each battle to be evaluated against the territorial
perimeter that most closely represents its chronological context.

The observed mean distance to the historical perimeter is compared
against simulated distributions.

Two alternative approaches are evaluated.

First, the settlement-based model is used to generate simulated battle
locations. Their distances to the corresponding historical perimeters
establish the expected boundary distances under a settlement conditioned
null model.

A complementary set of simulations uses exponential distance-decay
functions to generate spatial distributions increasingly concentrated
around territorial boundaries. These alternative distributions provide a
way to examine how different levels of border concentration compare with
the observed geographic pattern.

Second, boundary-weighted spatial distributions are constructed using
exponential distance decay:

{[} w(d) = e\^{}\{-d/\lambda\} {]}

where:

\begin{itemize}
\tightlist
\item
  (d) is the distance from a grid cell to the corresponding territorial
  perimeter.
\item
  (\lambda) is the distance-decay parameter.
\end{itemize}

Four decay distances are evaluated:

\begin{itemize}
\tightlist
\item
  25 kilometers
\item
  50 kilometers
\item
  100 kilometers
\item
  200 kilometers
\end{itemize}

Smaller decay parameters concentrate simulated locations more strongly
near the territorial perimeter.

The observed boundary distances are compared against the resulting
simulated distributions to evaluate whether the observed patterns are
consistent with different levels of territorial boundary association.

For the classification-based comparison, the border-weighted simulations
are also conducted separately for external battles and civil conflicts.
Each simulation preserves the observed number of battle-location records
within each classification and historical period. The period-specific
observed and simulated mean distances are then pooled across the full
historical dataset using the number of battle-location records in each
period as weights. This allows the overall results to account for
differences in the number of recorded battles across historical periods
while maintaining separate estimates for each conflict classification.

\subsection{Military and Civic
Infrastructure}\label{military-and-civic-infrastructure}

The final stage examines whether the geographic distribution of battles
is associated with specialized settlement infrastructure and whether
these relationships differ by conflict classification. The analysis
considers ports, temples, forts, and entertainment facilities, with
amphitheatres and theatres grouped into a single category. Two separate
approaches are used. The first evaluates whether battles are unusually
close to particular infrastructure features relative to the broader
settlement distribution. The second examines whether
infrastructure-weighted spatial distributions can reproduce the observed
clustering of external battles and civil conflicts.

The selection of these infrastructure categories is informed by the
preceding feature-proximity analysis, which identifies differences in
the relationships between conflict classifications and specialized
settlement features. Forts are used to represent military infrastructure
relevant to external battles, while entertainment facilities provide the
corresponding civic infrastructure category for civil conflicts.

\subsubsection{Proximity to Specialized Settlement
Features}\label{proximity-to-specialized-settlement-features}

The first analysis evaluates the proximity of battle locations to four
categories of settlement infrastructure: ports, temples, forts, and
entertainment facilities. For each category, the nearest-feature
distance is calculated for every observed battle location. The observed
mean distance is then compared with a conditional settlement-based Monte
Carlo model.

To prevent the feature being tested from directly influencing its own
background distribution, the corresponding feature category is excluded
from the settlement dataset before fitting the kernel density estimation
(KDE). A settlement intensity surface is then constructed from the
remaining settlement locations. This surface provides the sampling
probabilities used to generate simulated battle locations, representing
the expected distribution of battles based on the broader settlement
geography rather than the infrastructure category being tested.

The number of simulated locations matches the observed number of battle
locations in each classification. This process is repeated 999 times for
each feature category and classification, producing a simulated
distribution of mean distances to the nearest feature.

The observed mean distance is compared with the simulated distribution
using a lower-tail Monte Carlo test. This evaluates whether battles are
unusually close to a particular infrastructure category relative to what
would be expected from the broader settlement distribution. The
observed-to-expected ratio is calculated by dividing the observed mean
distance by the mean of the simulated distances. Ratios below 1 indicate
closer proximity than expected, while ratios above 1 indicate greater
distances than expected.

\subsubsection{Classification-Based Infrastructure
Models}\label{classification-based-infrastructure-models}

A separate simulation model examines whether the spatial clustering of
external battles and civil conflicts is consistent with distributions
weighted toward different types of infrastructure. External battles are
modeled using proximity to forts, while civil conflicts are modeled
using proximity to entertainment facilities. This approach differs from
the preceding analysis because it evaluates the clustering of battles
relative to one another, rather than their individual distances to
infrastructure features.

The analysis first consolidates multiple location records belonging to
the same battle into a single representative point. For battles with
multiple location records, the geometries are combined and their
centroid is used as the representative location. This produces a dataset
of 332 unique battles, consisting of 252 external battles and 80 civil
conflicts. The observed mean nearest-neighbor distance is then
calculated separately for external battles, civil conflicts, and the
combined dataset.

A shared 20 km sampling grid is constructed across the study area. For
each grid cell, the distance to the nearest relevant infrastructure
location is calculated. These distances are converted into sampling
weights using an exponential distance-decay function:

{[} w(d) = e\^{}\{-(d-d\_\{\min\})/\lambda\} {]}

where (d) represents the distance from a grid cell to the nearest
relevant infrastructure feature, (d\_\{\min\}) is the minimum distance
across the sampling grid, and (\lambda) is the distance-decay parameter.
Subtracting the minimum distance ensures that the highest sampling
weight is 1. Smaller decay parameters concentrate simulated locations
more strongly around the selected infrastructure, while larger
parameters produce broader spatial distributions.

Four decay parameters are evaluated: 25, 50, 100, and 200 km. For each
parameter, 999 simulations are generated using the corresponding
infrastructure-weighted sampling distribution. Each simulated
classification contains the same number of unique battles as its
observed group: 252 external battles and 80 civil conflicts. Simulated
locations are sampled from the weighted grid, with a small amount of
spatial jitter applied within each grid cell to avoid placing all
simulated points at identical cell locations.

The mean nearest-neighbor distance is calculated separately for each
simulated classification. The simulated external and civil-conflict
locations are also combined to calculate the corresponding statistic for
all battles. This combined result represents the clustering produced by
the two separately simulated classification groups rather than a single
distribution weighted toward one infrastructure category.

Finally, the infrastructure-weighted results are compared with a general
settlement-based baseline. This baseline uses the broader settlement
distribution to generate simulated battle locations without weighting
them toward a specific infrastructure category. Comparing the two models
provides a way to evaluate whether proximity to forts and entertainment
facilities offers a better explanation of the observed spatial
clustering than settlement geography alone.

\subsubsection{Representative Spatial
Simulations}\label{representative-spatial-simulations}

To illustrate the spatial patterns generated by the
infrastructure-weighted models, representative simulations are selected
from the 25 km decay model. For each classification, the simulation
whose mean nearest-neighbor distance is closest to the median of the 999
simulated values is selected. These simulations provide a representative
example of the spatial distributions generated by the model and are
compared with the observed battle locations.

\section{Results}\label{results}

\subsection{Spatial Randomness}\label{spatial-randomness}

The observed distribution of Roman battle locations exhibited
substantially greater spatial clustering than expected under complete
spatial randomness (CSR). Across 342 battle-location records, the
observed mean nearest-neighbor distance was 41.59 km, compared with an
expected mean distance of 108.70 km under the CSR model. The observed to
expected ratio was 0.383, indicating that the observed mean distance was
approximately 61.7\% lower than expected under spatial randomness.

None of the 999 simulated CSR distributions produced a mean
nearest-neighbor distance as extreme as the observed value. The
resulting Monte Carlo p-value was 0.001, providing strong evidence
against the hypothesis that battle locations were randomly distributed
across the study area.

Figure Figure~\ref{fig-spatial-randomness-distribution} shows the
distribution of simulated mean nearest-neighbor distances under CSR,
with the observed statistic indicated by a vertical reference line. The
position of the observed value relative to the simulated distribution
illustrates the degree to which battle locations were more spatially
concentrated than expected under a random spatial process.

\begin{figure}

\centering{

\includegraphics{figures/spatial_randomness_distribution.png}

}

\caption{\label{fig-spatial-randomness-distribution}Distribution of
simulated mean nearest-neighbor distances under complete spatial
randomness, with the observed statistic indicated by a vertical
reference line.}

\end{figure}%

\subsection{Settlement Based Models}\label{settlement-based-models}

Across the full historical dataset of battle-location records, the
observed mean nearest-neighbor distance was 41.59 km, compared with an
expected mean distance of 70.54 km under the settlement-based model. The
observed-to-expected ratio was 0.590, indicating that the observed mean
distance was approximately 41.0\% lower than expected based on
settlement geography an improvement over the complete spatial randomness
model, but spatial clustering of battle locations persisted even after
accounting for the distribution of known settlements.

\begin{figure}

\centering{

\includegraphics{figures/settlement_distribution.png}

}

\caption{\label{fig-settlement-distribution}Distribution of simulated
mean nearest-neighbor distances under the settlement-based model, with
the observed statistic indicated by a vertical reference line.}

\end{figure}%

\subsubsection{Border Distance
Investigation}\label{border-distance-investigation}

The border-distance investigation evaluated whether battle locations
were closer to historical Roman territorial boundaries than expected
based on the geographic distribution of known settlements. Across the
342 battle-location records, the observed mean distance to the nearest
border was 139.63 km, compared with an expected mean distance of 288.93
km under the settlement intensity model. The observed-to-expected ratio
was 0.483, indicating that battles occurred approximately 51.7\% closer
to territorial boundaries than expected based on settlement geography.
None of the 999 simulations produced a mean distance as small as the
observed value, suggesting that the relationship between battle
locations and Roman territorial boundaries warranted further
investigation beyond the influence of settlement distribution alone.

\begin{figure}

\centering{

\includegraphics{figures/border_distance_distribution.png}

}

\caption{\label{fig-border-distance-distribution}Distribution of
simulated mean distances to the nearest Roman territorial boundary under
the settlement-intensity model, with the observed mean distance
indicated by a vertical reference line.}

\end{figure}%

\subsection{Border Based Model}\label{border-based-model}

The simulated expected mean distance increased as the distance-decay
parameter became broader. Under the 25 km decay model, the expected mean
distance was 21.31 km, compared with 38.17 km under the 50 km model and
70.24 km under the 100 km model. Under the 200 km model, the expected
mean distance increased to 135.88 km, producing the closest
correspondence with the observed value. The observed-to-expected ratio
declined from 6.554 under the 25 km model to 1.028 under the 200 km
model.

Figure Figure~\ref{fig-border-distance-decay} compares the observed mean
border distance with the simulated distributions under each decay
parameter. The observed statistic exceeded the 97.5th percentile of the
simulations under the 25 km, 50 km, and 100 km models. Under the 200 km
model, it fell within the central 95\% simulation interval of 120.15 km
to 153.16 km. This indicates that the broader 200 km decay assumption
generated a distribution most consistent with the observed average
distance from battles to territorial boundaries.

In contrast to the settlement-based model, which found that battle
locations were more spatially clustered than expected given settlement
geography, the border model suggests that a relatively broad
distance-decay assumption is needed to reproduce the observed average
proximity to territorial boundaries.

This exploratory comparison does not establish that battles were caused
by, or deliberately positioned near, territorial boundaries. However,
the differences across decay parameters suggest that the relationship
between Roman territorial boundaries and battle locations warrants
further investigation.

\begin{figure}

\centering{

\includegraphics{figures/border_distance_decay.png}

}

\caption{\label{fig-border-distance-decay}Observed mean distance to the
nearest historical boundary compared with simulated distributions under
four distance-decay parameters.}

\end{figure}%

\subsubsection{Civil versus External
Conflict}\label{civil-versus-external-conflict}

The pooled results show a substantial difference between the two
conflict classifications. Across the 83 civil-conflict records, the
weighted observed mean distance to the nearest territorial boundary was
61.17 km. For the 259 external-conflict records, the corresponding
distance was 164.78 km. Thus, external battles occurred, on average,
approximately 2.7 times farther from territorial boundaries than civil
conflicts.

The observed-to-expected ratios also differed between classifications.
For civil conflicts, the ratio declined from 2.85 under the 25 km decay
model to 0.56 under the 200 km model. Under the 100 km model, the ratio
was 0.95, indicating that the observed mean distance was relatively
close to the simulated expectation. In contrast, external battles had
ratios above 1.0 under all four decay assumptions, ranging from 7.71
under the 25 km model to 1.14 under the 200 km model. This indicates
that external battles were consistently farther from territorial
boundaries than the border-weighted simulations predicted.

For all 342 battle-location records combined, the 200 km decay model
produced the closest correspondence between observed and expected mean
distances, with a ratio of 1.03. However, the classification-specific
results suggest that this overall correspondence masks differences
between external and civil conflicts. The 100 km and 200 km models
produced expected distances relatively close to those observed for civil
conflicts, while external battles remained farther from territorial
boundaries than expected under each model.

These results suggest that the spatial relationship between Roman
territorial boundaries and battle locations differed according to
conflict classification. Civil conflicts were generally located closer
to territorial boundaries, whereas external battles were more distant.
This comparison is descriptive and does not establish that the
differences are statistically significant; the period-level simulation
distributions provide additional context for examining how these
relationships varied over time.

\begin{longtable}[]{@{}
  >{\raggedright\arraybackslash}p{(\columnwidth - 12\tabcolsep) * \real{0.1111}}
  >{\raggedleft\arraybackslash}p{(\columnwidth - 12\tabcolsep) * \real{0.1481}}
  >{\raggedleft\arraybackslash}p{(\columnwidth - 12\tabcolsep) * \real{0.1481}}
  >{\raggedleft\arraybackslash}p{(\columnwidth - 12\tabcolsep) * \real{0.1481}}
  >{\raggedleft\arraybackslash}p{(\columnwidth - 12\tabcolsep) * \real{0.1481}}
  >{\raggedleft\arraybackslash}p{(\columnwidth - 12\tabcolsep) * \real{0.1481}}
  >{\raggedleft\arraybackslash}p{(\columnwidth - 12\tabcolsep) * \real{0.1481}}@{}}
\toprule\noalign{}
\begin{minipage}[b]{\linewidth}\raggedright
Classification
\end{minipage} & \begin{minipage}[b]{\linewidth}\raggedleft
Battle locations
\end{minipage} & \begin{minipage}[b]{\linewidth}\raggedleft
Observed mean (km)
\end{minipage} & \begin{minipage}[b]{\linewidth}\raggedleft
25 km ratio
\end{minipage} & \begin{minipage}[b]{\linewidth}\raggedleft
50 km ratio
\end{minipage} & \begin{minipage}[b]{\linewidth}\raggedleft
100 km ratio
\end{minipage} & \begin{minipage}[b]{\linewidth}\raggedleft
200 km ratio
\end{minipage} \\
\midrule\noalign{}
\endhead
\bottomrule\noalign{}
\endlastfoot
All Battles & 342 & 139.63 & 6.56 & 3.67 & 1.98 & 1.03 \\
Civil Conflict & 83 & 61.17 & 2.85 & 1.65 & 0.95 & 0.56 \\
External & 259 & 164.78 & 7.71 & 4.30 & 2.27 & 1.14 \\
\end{longtable}

\subsection{Settlement Attributes}\label{settlement-attributes}

\subsubsection{Battle Proximity to Specialized Settlement
Features}\label{battle-proximity-to-specialized-settlement-features}

To examine whether battle locations were unusually close to particular
types of settlement infrastructure, four feature categories were
evaluated: ports, temples, forts, and entertainment facilities
(amphitheatres and theatres). The analysis included 99 port locations,
54 temple locations, 50 fort locations, and 89 entertainment locations.
Results were evaluated separately for civil conflicts and external
battles.

\paragraph{Civil Conflicts}\label{civil-conflicts}

Civil conflicts exhibited the strongest proximity to entertainment
facilities. The observed mean distance to the nearest entertainment
location was 95.90 km, compared with a simulated expected distance of
229.26 km. The observed-to-expected ratio was 0.42, with a lower-tail
Monte Carlo p-value of 0.001. The observed distance was substantially
below the simulated distribution, indicating that civil conflicts were
unusually close to entertainment facilities relative to the broader
settlement geography.

Civil conflicts were also closer to ports than expected. The observed
mean distance was 134.03 km, compared with an expected distance of
183.53 km. This produced an observed-to-expected ratio of 0.73 and a
p-value of 0.004, providing evidence of unusually close proximity to
port locations.

The results for temples and forts did not show comparable evidence of
unusually close proximity. The observed mean distance to temples was
204.06 km, compared with an expected distance of 217.65 km (ratio =
0.94, p = 0.229). For forts, the observed mean distance was slightly
greater than expected, at 246.42 km compared with 240.74 km (ratio =
1.02, p = 0.634).

Overall, the civil-conflict results indicate that battles were
particularly close to entertainment facilities and ports, while the
results for temples and forts did not indicate unusually close
proximity.

\paragraph{External Battles}\label{external-battles}

External battles also exhibited unusually close proximity to
entertainment facilities. Their observed mean distance was 182.87 km,
compared with an expected distance of 229.57 km. The
observed-to-expected ratio was 0.80, with a lower-tail p-value of 0.002.

A similar pattern appeared for forts. External battles had an observed
mean distance of 210.15 km, compared with an expected distance of 241.12
km (ratio = 0.87, p = 0.005). This indicates that external battles were
closer to forts than expected under the settlement-based simulation.

The temple results were less conclusive. The observed mean distance was
200.74 km, compared with an expected distance of 217.03 km, producing a
ratio of 0.92 and a p-value of 0.051. Although the observed distance was
shorter than expected, the result was just above the conventional 0.05
threshold.

Ports showed a different pattern. External battles were farther from
ports than expected, with an observed mean distance of 218.88 km
compared with an expected distance of 184.55 km. The
observed-to-expected ratio was 1.19, and the lower-tail p-value was
1.000. Thus, the simulations provide no evidence that external battles
were unusually close to ports.

\paragraph{Comparison Between Conflict
Classifications}\label{comparison-between-conflict-classifications}

The results reveal differences in the spatial relationships between
battle classifications and specialized settlement features.
Entertainment facilities were the only category for which both civil
conflicts and external battles showed evidence of unusually close
proximity relative to the settlement-based simulations. However, the
difference between observed and expected distances was substantially
larger for civil conflicts.

Ports and forts displayed contrasting patterns across the two
classifications. Civil conflicts were significantly closer to ports than
expected, whereas external battles were farther away. Conversely,
external battles were significantly closer to forts than expected, while
civil conflicts showed no evidence of unusually close proximity to them.

Temples produced the least conclusive results. Neither classification
showed evidence of unusually close proximity at the 0.05 significance
threshold, although the external-battle result approached that
threshold.

Table Table~\ref{tbl-feature-proximity} summarizes the observed and
simulated distances for all eight comparisons. Figure
Figure~\ref{fig-feature-proximity} shows the simulated distributions
alongside the observed and expected mean distances.

\begin{longtable}[]{@{}
  >{\raggedright\arraybackslash}p{(\columnwidth - 14\tabcolsep) * \real{0.1250}}
  >{\raggedright\arraybackslash}p{(\columnwidth - 14\tabcolsep) * \real{0.1250}}
  >{\raggedleft\arraybackslash}p{(\columnwidth - 14\tabcolsep) * \real{0.1250}}
  >{\raggedleft\arraybackslash}p{(\columnwidth - 14\tabcolsep) * \real{0.1250}}
  >{\raggedleft\arraybackslash}p{(\columnwidth - 14\tabcolsep) * \real{0.1250}}
  >{\raggedleft\arraybackslash}p{(\columnwidth - 14\tabcolsep) * \real{0.1250}}
  >{\raggedleft\arraybackslash}p{(\columnwidth - 14\tabcolsep) * \real{0.1250}}
  >{\raggedleft\arraybackslash}p{(\columnwidth - 14\tabcolsep) * \real{0.1250}}@{}}
\caption{Observed and simulated mean distances between battle locations
and specialized settlement
features.}\label{tbl-feature-proximity}\tabularnewline
\toprule\noalign{}
\begin{minipage}[b]{\linewidth}\raggedright
Feature category
\end{minipage} & \begin{minipage}[b]{\linewidth}\raggedright
Classification
\end{minipage} & \begin{minipage}[b]{\linewidth}\raggedleft
Battle locations
\end{minipage} & \begin{minipage}[b]{\linewidth}\raggedleft
Feature locations
\end{minipage} & \begin{minipage}[b]{\linewidth}\raggedleft
Observed mean (km)
\end{minipage} & \begin{minipage}[b]{\linewidth}\raggedleft
Expected mean (km)
\end{minipage} & \begin{minipage}[b]{\linewidth}\raggedleft
Ratio
\end{minipage} & \begin{minipage}[b]{\linewidth}\raggedleft
Lower-tail p-value
\end{minipage} \\
\midrule\noalign{}
\endfirsthead
\toprule\noalign{}
\begin{minipage}[b]{\linewidth}\raggedright
Feature category
\end{minipage} & \begin{minipage}[b]{\linewidth}\raggedright
Classification
\end{minipage} & \begin{minipage}[b]{\linewidth}\raggedleft
Battle locations
\end{minipage} & \begin{minipage}[b]{\linewidth}\raggedleft
Feature locations
\end{minipage} & \begin{minipage}[b]{\linewidth}\raggedleft
Observed mean (km)
\end{minipage} & \begin{minipage}[b]{\linewidth}\raggedleft
Expected mean (km)
\end{minipage} & \begin{minipage}[b]{\linewidth}\raggedleft
Ratio
\end{minipage} & \begin{minipage}[b]{\linewidth}\raggedleft
Lower-tail p-value
\end{minipage} \\
\midrule\noalign{}
\endhead
\bottomrule\noalign{}
\endlastfoot
Entertainment & Civil Conflict & 83 & 89 & 95.90 & 229.26 & 0.42 &
0.001 \\
Ports & Civil Conflict & 83 & 99 & 134.03 & 183.53 & 0.73 & 0.004 \\
Temples & Civil Conflict & 83 & 54 & 204.06 & 217.65 & 0.94 & 0.229 \\
Forts & Civil Conflict & 83 & 50 & 246.42 & 240.74 & 1.02 & 0.634 \\
Entertainment & External & 259 & 89 & 182.87 & 229.57 & 0.80 & 0.002 \\
Forts & External & 259 & 50 & 210.15 & 241.12 & 0.87 & 0.005 \\
Temples & External & 259 & 54 & 200.74 & 217.03 & 0.92 & 0.051 \\
Ports & External & 259 & 99 & 218.88 & 184.55 & 1.19 & 1.000 \\
\end{longtable}

\begin{figure}

\centering{

\includegraphics{visualization/battle_feature_proximity.png}

}

\caption{\label{fig-feature-proximity}Simulated distributions of mean
distance to the nearest specialized settlement feature, separated by
battle classification. Vertical lines indicate observed and expected
mean distances.}

\end{figure}%

These findings suggest that the relationship between Roman battles and
settlement infrastructure varied by conflict classification and feature
type. Civil conflicts were particularly close to entertainment
facilities and ports, while external battles were closer than expected
to entertainment facilities and forts. These patterns are consistent
with different spatial relationships between the two types of conflict
and the surrounding settlement landscape, although they do not establish
that these features directly influenced where battles occurred.

The results should also be interpreted in light of the eight comparisons
conducted. The p-values are exploratory and have not been adjusted for
multiple testing. In addition, the simulations provide a conditional
comparison against the broader settlement distribution rather than a
uniformly random distribution of locations across the study area.

\subsection{Military and Civic
Infrastructure}\label{military-and-civic-infrastructure-1}

\subsubsection{Unique Battle Locations}\label{unique-battle-locations}

The observed mean nearest-neighbor distances were 50.25 km for external
battles, 89.65 km for civil conflicts, and 43.02 km for all battles
combined. These results indicate that external battles were more
spatially concentrated than civil conflicts, while the combined dataset
exhibited the shortest mean nearest-neighbor distance.

\subsubsection{General Settlement
Baseline}\label{general-settlement-baseline}

As established in the preceding settlement-based analysis (Figure
Figure~\ref{fig-settlement-distribution}), battle locations remained
more spatially clustered than expected under a model accounting for the
geographic distribution of known settlements.

The expected mean nearest-neighbor distances were 82.94 km for external
battles, 153.46 km for civil conflicts, and 71.87 km for all battles
combined. The observed-to-expected ratios were 0.606, 0.584, and 0.599,
respectively, with two-sided Monte Carlo p-values of 0.002 for all three
comparisons. These settlement-based results provide the baseline against
which the infrastructure-weighted simulations are compared below.

\subsubsection{Comparison of Settlement and Infrastructure
Models}\label{comparison-of-settlement-and-infrastructure-models}

The infrastructure-weighted simulations produced different results
depending on the decay distance and battle classification. In general,
increasing the decay distance produced less spatially concentrated
simulated patterns, as reflected by progressively larger expected mean
nearest-neighbor distances.

For external battles, the 25 km fort-based model produced an expected
mean nearest-neighbor distance of 39.14 km, compared with the observed
distance of 50.25 km. The observed-to-expected ratio was 1.284, and the
two-sided Monte Carlo p-value was 0.002. In this case, the observed
battles were less tightly clustered than the simulated locations
concentrated around forts.

As the decay distance increased, the expected nearest-neighbor distance
also increased. Under the 50 km model, the expected distance was 62.06
km (ratio = 0.810, p = 0.002). The 100 km and 200 km models produced
expected distances of 87.97 km and 108.65 km, respectively, with
observed-to-expected ratios of 0.571 and 0.463. Both models produced
p-values of 0.002. Thus, the broader fort-weighted distributions
produced simulated battle patterns that were less tightly clustered than
the observed external battles.

Civil conflicts displayed a different pattern. The 25 km
entertainment-based model produced an expected mean nearest-neighbor
distance of 111.21 km, compared with the observed distance of 89.65 km
(ratio = 0.806, p = 0.020). The 50 km model produced an expected
distance of 129.98 km (ratio = 0.690, p = 0.002). Under the 100 km and
200 km models, expected distances increased to 154.72 km and 186.27 km,
respectively. Both models produced observed-to-expected ratios below
0.60 and p-values of 0.002. Across all four decay parameters, the
simulated civil-conflict patterns were less tightly clustered than the
observed distribution.

For all battles combined, the 25 km infrastructure model produced an
expected mean nearest-neighbor distance of 45.23 km, relatively close to
the observed distance of 43.02 km. Its observed-to-expected ratio was
0.951, and the two-sided p-value was 0.304. This was the only model
comparison that did not indicate a statistically unusual difference
between the observed and simulated clustering statistics at the
conventional 0.05 threshold.

The broader infrastructure models produced progressively larger expected
distances. At 50 km, the expected distance was 62.39 km (ratio = 0.690,
p = 0.002). At 100 km and 200 km, expected distances increased to 80.89
km and 96.54 km, with ratios of 0.532 and 0.446, respectively. Both
comparisons produced p-values of 0.002. For comparison, the general
settlement model produced an expected distance of 71.87 km (ratio =
0.599, p = 0.002).

Overall, the results indicate that the relationship between
infrastructure-weighted spatial distributions and observed battle
clustering varied by classification. The 25 km fort-based model
generated a more tightly clustered external-battle pattern than the
observed distribution, while the broader fort-based models generated
less tightly clustered patterns. Civil-conflict simulations consistently
produced less clustered patterns than those observed, although the
difference was smallest under the 25 km entertainment-based model.

The 25 km infrastructure model also produced the closest correspondence
with the observed clustering of all battles combined. However, this
result should be interpreted in the context of the
classification-specific simulations: the combined model joins separately
simulated external battles weighted toward forts and civil conflicts
weighted toward entertainment facilities. It therefore represents a
combination of two infrastructure-weighted distributions rather than a
single probability surface.

Table Table~\ref{tbl-infrastructure-model} summarizes the observed and
simulated mean nearest-neighbor distances across the general settlement
and infrastructure-weighted models.

\begin{longtable}[]{@{}
  >{\raggedright\arraybackslash}p{(\columnwidth - 14\tabcolsep) * \real{0.1250}}
  >{\raggedright\arraybackslash}p{(\columnwidth - 14\tabcolsep) * \real{0.1250}}
  >{\raggedleft\arraybackslash}p{(\columnwidth - 14\tabcolsep) * \real{0.1250}}
  >{\raggedleft\arraybackslash}p{(\columnwidth - 14\tabcolsep) * \real{0.1250}}
  >{\raggedleft\arraybackslash}p{(\columnwidth - 14\tabcolsep) * \real{0.1250}}
  >{\raggedleft\arraybackslash}p{(\columnwidth - 14\tabcolsep) * \real{0.1250}}
  >{\raggedleft\arraybackslash}p{(\columnwidth - 14\tabcolsep) * \real{0.1250}}
  >{\raggedleft\arraybackslash}p{(\columnwidth - 14\tabcolsep) * \real{0.1250}}@{}}
\caption{Comparison of observed and simulated mean nearest-neighbor
distances across settlement and infrastructure-weighted
models.}\label{tbl-infrastructure-model}\tabularnewline
\toprule\noalign{}
\begin{minipage}[b]{\linewidth}\raggedright
Model
\end{minipage} & \begin{minipage}[b]{\linewidth}\raggedright
Classification
\end{minipage} & \begin{minipage}[b]{\linewidth}\raggedleft
Observed mean (km)
\end{minipage} & \begin{minipage}[b]{\linewidth}\raggedleft
Expected mean (km)
\end{minipage} & \begin{minipage}[b]{\linewidth}\raggedleft
Ratio
\end{minipage} & \begin{minipage}[b]{\linewidth}\raggedleft
2.5th percentile (km)
\end{minipage} & \begin{minipage}[b]{\linewidth}\raggedleft
97.5th percentile (km)
\end{minipage} & \begin{minipage}[b]{\linewidth}\raggedleft
Two-sided p-value
\end{minipage} \\
\midrule\noalign{}
\endfirsthead
\toprule\noalign{}
\begin{minipage}[b]{\linewidth}\raggedright
Model
\end{minipage} & \begin{minipage}[b]{\linewidth}\raggedright
Classification
\end{minipage} & \begin{minipage}[b]{\linewidth}\raggedleft
Observed mean (km)
\end{minipage} & \begin{minipage}[b]{\linewidth}\raggedleft
Expected mean (km)
\end{minipage} & \begin{minipage}[b]{\linewidth}\raggedleft
Ratio
\end{minipage} & \begin{minipage}[b]{\linewidth}\raggedleft
2.5th percentile (km)
\end{minipage} & \begin{minipage}[b]{\linewidth}\raggedleft
97.5th percentile (km)
\end{minipage} & \begin{minipage}[b]{\linewidth}\raggedleft
Two-sided p-value
\end{minipage} \\
\midrule\noalign{}
\endhead
\bottomrule\noalign{}
\endlastfoot
Settlement & External & 50.25 & 82.94 & 0.61 & 75.14 & 90.61 & 0.002 \\
Infrastructure 25 km & External & 50.25 & 39.14 & 1.28 & 35.23 & 43.12 &
0.002 \\
Infrastructure 50 km & External & 50.25 & 62.06 & 0.81 & 56.49 & 68.17 &
0.002 \\
Infrastructure 100 km & External & 50.25 & 87.97 & 0.57 & 80.84 & 95.14
& 0.002 \\
Infrastructure 200 km & External & 50.25 & 108.65 & 0.46 & 100.43 &
116.86 & 0.002 \\
Settlement & Civil Conflict & 89.65 & 153.46 & 0.58 & 127.27 & 178.88 &
0.002 \\
Infrastructure 25 km & Civil Conflict & 89.65 & 111.21 & 0.81 & 93.36 &
129.79 & 0.020 \\
Infrastructure 50 km & Civil Conflict & 89.65 & 129.98 & 0.69 & 111.03 &
148.92 & 0.002 \\
Infrastructure 100 km & Civil Conflict & 89.65 & 154.72 & 0.58 & 133.54
& 175.50 & 0.002 \\
Infrastructure 200 km & Civil Conflict & 89.65 & 186.27 & 0.48 & 161.90
& 213.45 & 0.002 \\
Settlement & All Battles & 43.02 & 71.87 & 0.60 & 65.76 & 77.63 &
0.002 \\
Infrastructure 25 km & All Battles & 43.02 & 45.23 & 0.95 & 41.28 &
49.23 & 0.304 \\
Infrastructure 50 km & All Battles & 43.02 & 62.39 & 0.69 & 57.34 &
67.31 & 0.002 \\
Infrastructure 100 km & All Battles & 43.02 & 80.89 & 0.53 & 75.18 &
86.98 & 0.002 \\
Infrastructure 200 km & All Battles & 43.02 & 96.54 & 0.45 & 89.81 &
102.66 & 0.002 \\
\end{longtable}

Figure Figure~\ref{fig-infrastructure-distributions} shows the
distributions of simulated mean nearest-neighbor distances for the
general settlement model and the four infrastructure-weighted models.
The observed mean nearest-neighbor distance is shown for each
classification. The distributions illustrate how changing the
infrastructure distance-decay parameter affects the expected spatial
clustering of simulated battles.

\begin{figure}

\centering{

\includegraphics{figures/infrastructure_simulation_distributions.png}

}

\caption{\label{fig-infrastructure-distributions}Simulated mean
nearest-neighbor distance distributions for the general settlement model
and four infrastructure-weighted models, separated by battle
classification.}

\end{figure}%

\subsubsection{Representative Spatial
Simulations}\label{representative-spatial-simulations-1}

Figure Figure~\ref{fig-representative-infrastructure} compares the
observed battle locations with the representative 25 km
infrastructure-weighted simulations for external battles and civil
conflicts. This comparison provides a geographic complement to the
summary nearest-neighbor statistics, allowing the spatial configurations
produced by the model to be examined alongside the observed
distribution.

\begin{figure}

\centering{

\includegraphics{figures/representative_infrastructure_simulations.png}

}

\caption{\label{fig-representative-infrastructure}Observed battle
locations and representative 25 km infrastructure-weighted simulations
for external battles and civil conflicts.}

\end{figure}%

The simulation results demonstrate that infrastructure-weighted
distributions can reproduce some aspects of the observed spatial
clustering, but their correspondence varies by conflict classification
and distance-decay assumption. The results do not establish that battles
were caused by, or deliberately positioned near, forts or entertainment
facilities. Rather, they suggest that the spatial relationship between
battle locations and different forms of infrastructure merits further
investigation.

\section{Conclusion}\label{conclusion}

The spatial analyses identify several distinct patterns in the
geographic distribution of Roman battles. Battle locations exhibit
substantial clustering relative to a uniform random distribution, and
this clustering remains apparent when settlement geography is
incorporated as a base model. Battles are also significantly closer to
historical Roman territorial boundaries than expected under the
settlement based model. Finally, the infrastructure analysis suggests
that the relationship between military activity and specialized
settlement features differs between external and civil conflicts.

\subsection{Interpretation}\label{interpretation}

Taken together, these simulations suggest that the general settlement
distribution alone does not reproduce the observed degree of battle
clustering. The infrastructure-weighted models provide a more specific
comparison, showing that the relationship between infrastructure
proximity and battle clustering depends on both the type of conflict and
the assumed distance-decay parameter.

The 25 km infrastructure model produced the closest overall
correspondence between observed and simulated clustering, particularly
when all battles were combined. However, this aggregate result conceals
differences between the two classifications: the observed external
battles were less clustered than the fort-weighted simulation, while
civil conflicts were more clustered than the entertainment-weighted
simulation.

These findings provide evidence that settlement infrastructure may help
characterize the spatial distribution of Roman battles, but they do not
establish that forts or entertainment facilities directly determined
battle locations. The results are conditional on the selected features,
KDE and sampling-grid parameters, and the assumed distance-decay
functions. The comparisons are also exploratory and involve multiple
model tests.




\end{document}
